\section{Introduction}
\label{sec:intro}

The rapid adoption of large language models as agents~\citep{tot,react,selfconsistency} has turned inference from a single continuation into a search process.
An agentic step expands a committed trunk of length $L$ into $k$ residuals and later keeps a single spine (Figure~\ref{fig:problem}).
This pattern is the serving unit of Tree-of-Thoughts, self-consistency, and tool-using loops, and it is where latency and memory are spent.
However, existing stacks exhibit a critical inefficiency: residuals that the search will abort still incur GPU prefill, KV transfer, and decode until a prune score exists.
That waste grows with fan-out width and is paid on every step, posing a significant cost to the reliability of agentic inference in real-world applications.

The susceptibility of LLM serving to this waste stems from a fundamental limitation in parsing the fan-out, i.e., the inability to distinguish a residual that must be encoded from one that can be refused as text~\citep{esc,specrej,dpts}.
In reducing that cost, replacing the search algorithm, or retraining the model to emit fewer thoughts, can be computationally prohibitive and is orthogonal to serving.
Alternatively, existing mitigation strategies focus on scoring tokens that have already been generated~\citep{esc,specrej,dpts}, on recording a prefix only after a full block is hashed~\citep{vllm,sglang}, or on transferring every page a disaggregated prefiller has computed~\citep{distserve,splitwise}.
Such strategies often provide only decode-side savings, or they ship KV that the search is about to drop.
Additionally, imposing a single threshold on both prefill and decode can drop residuals that still solve, thereby undermining the quality of the kept spine.

\begin{figure*}[t]
\centering
\resizebox{\textwidth}{!}{\begin{tikzpicture}[
  x=1cm, y=1cm,
  font=\sffamily\scriptsize,
  >=Latex,
  line cap=round, line join=round,
  panel/.style={draw=fsblue!35, line width=0.55pt, rounded corners=2pt, fill=white},
  head/.style={font=\scriptsize\bfseries\sffamily, text=white, anchor=west},
  chip/.style={draw, rounded corners=1pt, inner sep=1.4pt, font=\tiny\sffamily, line width=0.4pt},
  arr/.style={-{Latex[length=1.6mm,width=1.2mm]}, line width=0.5pt, draw=fsblue!80},
  thought/.style={draw, rounded corners=1pt, minimum width=0.72cm, minimum height=0.38cm, inner sep=0pt, font=\tiny\sffamily},
]
\path (0,0) rectangle (16.35,6.55);

% ---------- panels ----------
\draw[panel] (0.00,0.62) rectangle (5.28,6.50);
\draw[panel] (5.48,0.62) rectangle (10.72,6.50);
\draw[panel] (10.92,0.62) rectangle (16.35,6.50);
\fill[fsblue] (0.00,6.08) rectangle (5.28,6.50);
\fill[fsorange!88!black] (5.48,6.08) rectangle (10.72,6.50);
\fill[fspurple!80!black] (10.92,6.08) rectangle (16.35,6.50);
\node[head] at (0.12,6.29) {(a) Search scores a written thought};
\node[head] at (5.60,6.29) {(b) Hosts share one prefill};
\node[head] at (11.04,6.29) {(c) Cache and transfer, after the page};

% ---------- (a) ToT-style tree, original redraw ----------
\node[draw=fsblue!60, fill=fslight, ellipse, inner sep=1.5pt, font=\tiny\sffamily] (inp) at (2.64,5.62) {Input};
\node[thought, fill=fscommit, draw=fsgreen!70!black] (t0) at (1.15,4.55) {};
\node[thought, fill=fsgreen!55!black, draw=fsgreen!70!black] (t1) at (2.64,4.55) {};
\node[thought, fill=fsmiss, draw=fsred!70] (t2) at (4.13,4.55) {};
\node[thought, fill=fscommit, draw=fsgreen!70!black] (t00) at (0.72,3.55) {};
\node[thought, fill=fsmiss, draw=fsred!70] (t01) at (1.58,3.55) {};
\node[thought, fill=fscommit, draw=fsgreen!70!black] (t10) at (2.28,3.55) {};
\node[thought, fill=fsmiss, draw=fsred!70] (t11) at (3.10,3.55) {};
\node[draw=fsgreen!60!black, fill=fscommit, ellipse, inner sep=1.4pt, font=\tiny\sffamily] (out) at (1.50,2.55) {Output};
\draw[arr, draw=fsgreen!70!black] (inp) -- (t0);
\draw[arr, draw=fsgreen!70!black] (inp) -- (t1);
\draw[arr, draw=fsred!70] (inp) -- (t2);
\draw[arr, draw=fsgreen!70!black] (t0) -- (t00);
\draw[arr, draw=fsred!70] (t0) -- (t01);
\draw[arr, draw=fsgreen!70!black] (t1) -- (t10);
\draw[arr, draw=fsred!70] (t1) -- (t11);
\draw[arr, draw=fsgreen!70!black] (t00) -- (out);
\node[font=\tiny\sffamily, text=fsgreen!60!black, anchor=west] at (3.35,2.62) {kept};
\node[font=\tiny\sffamily, text=fsred, anchor=west] at (3.35,2.28) {dropped};
\fill[fscommit] (3.15,2.52) rectangle (3.30,2.70);
\draw[fsgreen!70!black] (3.15,2.52) rectangle (3.30,2.70);
\fill[fsmiss] (3.15,2.18) rectangle (3.30,2.36);
\draw[fsred!70] (3.15,2.18) rectangle (3.30,2.36);

\node[chip, draw=fsorange!70, fill=fsspec, anchor=west] (prop) at (0.18,1.55) {propose};
\node[chip, draw=fsblue!60, fill=fslight, anchor=west] (val) at (2.55,1.55) {value: sure / maybe / impossible};
\draw[arr] (prop.east) -- (val.west);
\node[font=\tiny\sffamily, text=fsgray, align=left, anchor=west] at (0.18,0.95)
  {The value prompt reads a thought\\that the propose step has already written.};

% ---------- (b) three hosts, identical prefill ----------
\foreach \y/\name/\dw/\cut in {
  4.95/ESC/2.05/window,
  3.85/SR/1.35/{$\tau{=}256$},
  2.75/DPTS/0.85/{step 100}
} {
  \node[font=\tiny\sffamily\bfseries, anchor=west] at (5.62,\y) {\name};
  \fill[fsgray!35] (6.55,\y-0.16) rectangle (7.85,\y+0.16);
  \draw[fsgray] (6.55,\y-0.16) rectangle (7.85,\y+0.16);
  \fill[fsblue!45] (7.85,\y-0.16) rectangle ({7.85+\dw},\y+0.16);
  \draw[fsblue!80] (7.85,\y-0.16) rectangle ({7.85+\dw},\y+0.16);
  \draw[fsred, line width=0.9pt] ({7.85+\dw},\y-0.28) -- ({7.85+\dw},\y+0.28);
  \node[font=\tiny\sffamily, text=fsred, anchor=west] at ({7.95+\dw},\y) {\cut};
}
\node[font=\tiny\sffamily, text=fsgray, anchor=west] at (6.55,2.22) {prefill};
\node[font=\tiny\sffamily, text=fsblue, anchor=west] at (7.95,2.22) {decode until the score};
\node[font=\tiny\sffamily, align=left, anchor=west] at (5.62,1.35)
  {$k{=}16$, one GSM8K run:\\prefill tokens $= 5008$ for all three.\\They differ only after the kernel.};
\node[font=\tiny\sffamily, text=fsgray, align=left, anchor=west] at (5.62,0.78)
  {ESC keeps the window open.\\SR and DPTS still decode to the cut.};

% ---------- (c) vLLM + DistServe, original redraw ----------
\node[chip, draw=fsgreen!60!black, fill=fscommit] (sch) at (11.70,5.55) {Scheduler};
\node[chip, draw=fsblue!50, fill=fslight, align=center] (kvm) at (13.55,5.55) {KV manager\\[-1pt]\tiny block table};
\node[chip, draw=fsblue!50, fill=white] (wk) at (15.45,5.55) {Worker};
\draw[arr] (sch) -- (kvm);
\draw[arr] (kvm) -- (wk);
\node[font=\tiny\sffamily, text=fsgray, anchor=west] at (11.10,4.85)
  {Hash after a full block.};

\node[chip, draw=fsred!60, fill=fsmiss] (ctl) at (11.35,4.15) {Controller};
\node[draw=fsorange!70, fill=fsspec, rounded corners=1.5pt, align=center,
      minimum width=1.55cm, minimum height=1.05cm, font=\tiny\sffamily] (pre) at (12.55,2.95) {};
\node[font=\tiny\sffamily\bfseries, text=fsorange!80!black] at (12.55,3.22) {Prefill};
\node[font=\tiny\sffamily] at (12.55,2.90) {instance};
\node[font=\tiny\sffamily, text=fsgray] at (12.55,2.60) {computed KV};
\node[draw=fsblue!60, fill=fslight, rounded corners=1.5pt, align=center,
      minimum width=1.55cm, minimum height=1.05cm, font=\tiny\sffamily] (dec) at (15.35,2.95) {};
\node[font=\tiny\sffamily\bfseries, text=fsblue] at (15.35,3.22) {Decode};
\node[font=\tiny\sffamily] at (15.35,2.90) {instance};
\node[font=\tiny\sffamily, text=fsgray] at (15.35,2.60) {sampler};
\draw[arr] (ctl.south) -- (pre.north);
\draw[arr, draw=fsred, line width=0.7pt] (pre.east) -- (dec.west);
\node[font=\tiny\sffamily, text=fsred, align=center] at (13.95,3.70) {no admission};
\node[font=\tiny\sffamily, text=fsgray, align=left, anchor=west] at (11.08,1.55)
  {The first fan-out has no page yet,\\so every residual is computed.\\An aborted residual still crosses.};

% ---------- bottom cost bar ----------
\fill[black!4] (0.00,0.00) rectangle (16.35,0.48);
\draw[fsgray!50] (0.00,0.48) -- (16.35,0.48);
\node[font=\tiny\sffamily, anchor=west] at (0.15,0.24)
  {\textbf{Paid before abort:}\quad
   $W$ GPU prefill \quad\textperiodcentered\quad
   $X$ connector payload \quad\textperiodcentered\quad
   decode-to-cut (window, $\tau$, or mini-step)};
\end{tikzpicture}
}
\caption{Illustration of wasted work in an agentic fan-out.
(a)~Tree-of-Thoughts writes a thought and then values it~\citep{tot}.
(b)~ESC, speculative rejection, and DPTS read a prefix that is already resident~\citep{esc,specrej,dpts}; on this GSM8K run their prefill token count is the same.
(c)~Automatic prefix caching hashes a full block~\citep{vllm}, and a disaggregated runtime transfers the KV the prefiller has already computed~\citep{distserve}.
GPU prefill $W$, connector payload $X$, and decode-to-cut are incurred before the abort.}
\label{fig:problem}
\end{figure*}

In this paper, we focus on the source of that inefficiency: the stack's confusion between a residual that must be prefilled and a residual that can be skipped.
We start from a fundamental question: \emph{what makes the difference between a residual that should be encoded and one that should not, from the serving stack's perspective?}
The engine relies on implicit criteria that only become observable after tokens exist.
Concretely, from a cost perspective, a residual is \emph{live} if some continuation of it can still produce a correct answer, and \emph{abortable} if it is a repeated loop or illegal text that the search will drop.
Waste arises when those implicit criteria fail to align with the live--abortable separation that is already present in the residual text.
Our approach is motivated by solving that misalignment with the following two general steps.

\begin{itemize}
  \item \textbf{Prefill scoring.}
    We identify abortable residuals by a score on the residual text, whose activations bias the stack toward encoding the same span as a thought rather than as skippable text.
  \item \textbf{Decode intervention.}
    We prune the identified residuals over the prefill of the fan-out, and we continue an uncertain residual only when a short generated prefix clears a second bar.
\end{itemize}

Specifically, we introduce two unshared scores with ordered bars $(\tau_{\mathrm{pre}}, \tau_{\mathrm{dec}})$, $\tau_{\mathrm{pre}} < \tau_{\mathrm{dec}}$.
The prefill score reads the residual and drops a repeated loop before the kernel and before any KV transfer.
The decode score reads a generated prefix after a probe of $t_0$ tokens and continues to the budget $D$ only when that prefix clears $\tau_{\mathrm{dec}}$.
We impose the invariant that index $0$ is always admitted, so pruning cannot erase the residual the search will keep if nothing else survives.
We show that the proposed split is host-invariant in prefill work: any decode-time pruner that scores only generated tokens shares one $W$, and that $W$ is exactly the prefill of the admitted set.

For compatibility with prefix caching and prefill--decode disaggregation~\citep{vllm,distserve,splitwise}, we apply pruning when encoding the KV cache of the residual, i.e., enabling efficient prompting when several residuals share the same trunk.
We accordingly refer to our approach as \sys.
Notably, \sys leaves the host's decode cut unchanged and is lightweight in computation, relying solely on a pair of string scores without extra test-time model calls.

Our contributions are summarized as follows:
\begin{itemize}
  \item We propose \sys, which reduces wasted prefill in agentic inference by identifying and pruning abortable residuals during KV-cache encoding of the residual.
    It steers the stack toward treating a loop as skippable text, while not interfering with the model's ability to follow live strategies.
  \item In identifying these residuals, we propose a split scoring mechanism that allows effective pruning from surface tests on the residual and on a short prefix.
    We also leverage an observed score gap---illegal text lies above $\tau_{\mathrm{pre}}$ on the residual and below $\tau_{\mathrm{dec}}$ on a still-illegal prefix---that further improves the fidelity of the split.
  \item In experiments, \sys significantly reduces the prefill shared by ESC, speculative rejection, and DPTS while preserving any-finished accuracy, and a short probe restores full MATH-500 accuracy at lower end-to-end time.
\end{itemize}
